\documentclass[../main.tex]{subfiles}

\begin{document}
\chapter{Desarrollo de plugins}

\section{Introducción}

SamaruC Editor incluye un sistema de plugins basado en \textbf{Java ServiceLoader} que permite extender la aplicación sin tocar su código fuente. Este documento explica paso a paso cómo crear, empaquetar e instalar un plugin.


\section{Arquitectura general}

La arquitectura de plugins se basa en una interfaz común \texttt{Plugin} que los desarrolladores implementan. SamaruC descubre e instancia estas clases automáticamente al arrancar, llamando a sus métodos de ciclo de vida para integrarlas en la aplicación, en la siguiente figura podemos ver un diagrama de la arquitectura del sistema de plugins.

\begin{figure}[!htb]
    \centering
    \includegraphics[width=14cm]{Diagrama_arquitectura_plugins.jpg}
    \caption{Diagrama de arquitectura del sistema de plugins.}
    \label{fig:Diagrama_arquitectura_plugins}
\end{figure}

\clearpage

El ciclo de vida de un plugin se describe en la siguiente tabla:

\begin{table}[ht]
    \centering
    \small
    \begin{tabular}{p{4cm}p{8cm}}
    \toprule
    \textbf{Método} & \textbf{Cuándo se llama} \\
    \midrule
    \texttt{onLoad(ctx)} & Justo después de cargar el JAR. Guarda el contexto aquí. \\
    \texttt{onEnable()} & Cuando la UI está lista. Agrega tus menús/componentes aquí. \\
    \texttt{onDisable()} & Cuando el plugin se desactiva o la app cierra. Limpia recursos aquí. \\
    \texttt{getName()} & En cualquier momento. Devuelve el nombre visible del plugin. \\
    \bottomrule
    \end{tabular}
    \caption{Ciclo de vida de un plugin}
    \label{tab:ciclo_vida_plugin}
\end{table}

\section{Requisitos previos}

\begin{itemize}
\item \textbf{JDK 17} o superior
\item \textbf{Maven 3.8+} (o Gradle si lo prefieres)
\item El JAR principal de SamaruC (\texttt{retro-editor-X.Y-SNAPSHOT.jar}) para compilar contra sus interfaces
\end{itemize}

\textbf{Consejo:} Añade el JAR de SamaruC a tu repositorio Maven local con:
\begin{lstlisting}[language=bash]
mvn install:install-file -Dfile=retro-editor-1.0-SNAPSHOT.jar \
  -DgroupId=com.retroeditor -DartifactId=retro-editor \
  -Dversion=1.0-SNAPSHOT -Dpackaging=jar
\end{lstlisting}

\clearpage

\section{Crear el proyecto Maven del plugin}

Crea la siguiente estructura de directorios:

\begin{figure}[!htb]
    \centering
    \includegraphics[width=14cm]{estructura_directorios.jpg}
    \caption{Estructura de directorios del proyecto Maven del plugin.}
    \label{fig:estructura_directorios_plugin}
\end{figure}

\texttt{pom.xml} mínimo:

\begin{lstlisting}[language=xml, caption=Configuración Maven básica]
<?xml version="1.0" encoding="UTF-8"?>
<project xmlns="http://maven.apache.org/POM/4.0.0"
         xmlns:xsi="http://www.w3.org/2001/XMLSchema-instance"
         xsi:schemaLocation="http://maven.apache.org/POM/4.0.0
             http://maven.apache.org/xsd/maven-4.0.0.xsd">
    <modelVersion>4.0.0</modelVersion>

    <groupId>com.miplugin</groupId>
    <artifactId>mi-plugin</artifactId>
    <version>1.0.0</version>
    <packaging>jar</packaging>

    <dependencies>
        <!-- Interfaces de SamaruC (solo compile-time) -->
        <dependency>
            <groupId>com.retroeditor</groupId>
            <artifactId>retro-editor</artifactId>
            <version>1.0-SNAPSHOT</version>
            <scope>provided</scope>
        </dependency>

        <!-- JavaFX (solo compile-time, lo provee el host) -->
        <dependency>
            <groupId>org.openjfx</groupId>
            <artifactId>javafx-controls</artifactId>
            <version>17.0.10</version>
            <scope>provided</scope>
        </dependency>
    </dependencies>

    <build>
        <plugins>
            <plugin>
                <groupId>org.apache.maven.plugins</groupId>
                <artifactId>maven-compiler-plugin</artifactId>
                <version>3.8.1</version>
                <configuration>
                    <source>17</source>
                    <target>17</target>
                </configuration>
            </plugin>
        </plugins>
    </build>
</project>
\end{lstlisting}

\textbf{Importante:} Las dependencias de SamaruC y JavaFX deben declararse con \texttt{<scope>provided</scope>} porque ya están disponibles en el classpath del host en tiempo de ejecución.

\section{Implementar la interfaz \texttt{Plugin}}

Crea \texttt{MiPlugin.java} implementando \texttt{com.retroeditor.plugin.Plugin}:

\begin{lstlisting}[language=java, caption=Implementación básica de un plugin]
package com.miplugin;

import com.retroeditor.plugin.Plugin;
import com.retroeditor.plugin.PluginContext;
import javafx.application.Platform;
import javafx.scene.control.Menu;
import javafx.scene.control.MenuItem;

public class MiPlugin implements Plugin {

    private PluginContext context;
    private Menu          miMenu;

    /** Se llama al cargar el JAR. Guarda el contexto para usarlo despues. */
    @Override
    public void onLoad(PluginContext context) {
        this.context = context;
        context.log("MiPlugin cargado correctamente.");
    }

    /** Se llama cuando la UI esta lista. Aqui agregamos nuestros elementos. */
    @Override
    public void onEnable() {
        // Toda modificacion de la UI debe hacerse en el hilo de JavaFX
        Platform.runLater(() -> {
            miMenu = new Menu("Mi Plugin");

            MenuItem saludar = new MenuItem("Saludar");
            saludar.setOnAction(e -> context.log("Hola desde MiPlugin!"));

            miMenu.getItems().add(saludar);

            // Registrar el menu: SamaruC lo eliminara automaticamente 
            // al desactivar el plugin
            context.registerMenu(miMenu);
        });
    }

    /** Se llama al deshabilitar o cerrar la app. Limpia tus recursos. */
    @Override
    public void onDisable() {
        // Los menus registrados con registerMenu() se limpian automaticamente,
        // pero aqui puedes liberar cualquier otro recurso propio.
        context.log("MiPlugin desactivado.");
    }

    /** Nombre visible en la pantalla de configuracion de SamaruC. */
    @Override
    public String getName() {
        return "Mi Plugin de Ejemplo";
    }
}
\end{lstlisting}

\section{Registrar el plugin con ServiceLoader}

Java ServiceLoader requiere un fichero descriptor en \texttt{META-INF/services/}.

Crea el archivo:
\texttt{src/main/resources/META-INF/services/com.retroeditor.plugin.Plugin}

Con el contenido (nombre completo de tu clase implementadora):

\begin{lstlisting}
com.miplugin.MiPlugin
\end{lstlisting}

Si tu plugin tiene \textbf{varias implementaciones}, ponlas una por línea:

\begin{lstlisting}
com.miplugin.MiPlugin
com.miplugin.OtroPlugin
\end{lstlisting}

\clearpage

\section{Usar el \texttt{PluginContext}}

La interfaz \texttt{PluginContext} proporciona los siguientes métodos:

\begin{table}[ht]
    \centering
    \small
    \begin{tabular}{p{4cm}p{8cm}}
    \toprule
    \textbf{Método} & \textbf{Descripción} \\
    \midrule
    \texttt{getPrimaryStage()} & Devuelve el \texttt{Stage} principal de la aplicación. Útil para crear diálogos modales. \\
    \texttt{getScene()} & Devuelve la \texttt{Scene} principal. Permite acceder al \texttt{root} para añadir nodos. \\
    \texttt{getMenuBar()} & Devuelve la \texttt{MenuBar} de la app. Puedes añadir menús directamente o usar \texttt{registerMenu()}. \\
    \texttt{log(String)} & Escribe un mensaje en la pestaña \textbf{Monitor} de SamaruC. \\
    \texttt{registerMenu(Menu)} & Añade un \texttt{Menu} a la barra de menú y registra la referencia para limpieza automática. \\
    \texttt{unregisterMenu(Menu)} & Elimina manualmente un menú registrado. \\
    \texttt{unregisterAll()} & Elimina todos los elementos registrados por este plugin. Lo llama SamaruC automáticamente. \\
    \bottomrule
    \end{tabular}
    \caption{Métodos de PluginContext}
    \label{tab:metodos_plugin_context}
\end{table}

\subsection{Ejemplo: abrir un diálogo con el Stage del host}

\begin{lstlisting}[language=java]
@Override
public void onEnable() {
    Platform.runLater(() -> {
        Menu menu = new Menu("Herramientas");
        MenuItem abrirDialogo = new MenuItem("Abrir mi dialogo");
        abrirDialogo.setOnAction(e -> {
            Stage dialogo = new Stage();
            dialogo.initOwner(context.getPrimaryStage());   // modal respecto al host
            dialogo.setTitle("Mi herramienta");
            // ... configurar la escena del dialogo ...
            dialogo.show();
        });
        menu.getItems().add(abrirDialogo);
        context.registerMenu(menu);
    });
}
\end{lstlisting}

\subsection{Ejemplo: escribir en el Monitor}

\begin{lstlisting}[language=java]
context.log("Procesando fichero: " + archivo.getName());
\end{lstlisting}

\section{Empaquetar el plugin como JAR}

Desde la raíz de tu proyecto de plugin, ejecuta:

\begin{lstlisting}[language=bash]
mvn package
\end{lstlisting}

Esto generará \texttt{target/mi-plugin-1.0.0.jar}. Ese es el fichero que debes instalar en SamaruC.

\textbf{Nota:} No uses el plugin \texttt{maven-shade} en el JAR del plugin para no duplicar clases del host (JavaFX, etc.). El JAR debe contener \textbf{solo} las clases propias del plugin.

\section{Instalar el plugin en SamaruC}

Tienes dos formas de instalar un plugin:

\subsection{Opción A — Desde la interfaz de SamaruC (recomendado)}

\begin{enumerate}
\item Abre SamaruC Editor.
\item Ve a \textbf{Archivo / Configuración} (o el atajo de teclado correspondiente).
\item Selecciona la pesta\~na \textbf{Plugins}.
\item Haz clic en \textbf{Instalar plugin} y selecciona tu fichero \texttt{.jar}.
\item SamaruC copiara el JAR a la carpeta \texttt{plugins/} junto al ejecutable.
\item Pulsa \textbf{Recargar plugins} para activarlo sin reiniciar.
\end{enumerate}

\clearpage

\subsection{Opción B — Copiado manual}

Copia el fichero \texttt{.jar} directamente en la carpeta \texttt{plugins/} que se encuentra en el directorio de trabajo de SamaruC (al lado del ejecutable \texttt{RetroEditor.exe} o del JAR principal):

\begin{lstlisting}[language=bash]
RetroEditor/
+-- RetroEditor.exe
+-- app/
|   +-- retro-editor-1.0-SNAPSHOT.jar
+-- plugins/              <- aqui va tu JAR
    +-- mi-plugin-1.0.0.jar
\end{lstlisting}

Reinicia la aplicación y el plugin se cargará automáticamente.

\section{Habilitar y deshabilitar plugins}

En la pestaña \textbf{Plugins} de la pantalla de configuración puedes:

\begin{itemize}
    \item Ver la lista de todos los JARs descubiertos en \texttt{plugins/}.
    \item Ver su estado (habilitado / deshabilitado).
    \item Pulsar \textbf{Habilitar} o \textbf{Deshabilitar} para cambiar el estado.
    \item Pulsar \textbf{Recargar plugins} para aplicar cambios sin reiniciar.
\end{itemize}

El estado de cada plugin se guarda en el fichero:

\begin{lstlisting}
~/.retroeditor.plugins.properties
\end{lstlisting}

(donde \textasciitilde{} es tu carpeta de usuario del sistema operativo)

\clearpage

\section{Ejemplo completo}

A continuación se muestra un plugin mínimo funcional que añade un menú \textbf{``Utilidades''} con una opción que muestra la ruta del archivo actualmente abierto.

\subsection{Estructura del proyecto}

\begin{lstlisting}[language=bash]
plugin-utilidades/
+-- pom.xml
+-- src/main/
    +-- java/com/ejemplo/UtilsPlugin.java
    +-- resources/META-INF/services/com.retroeditor.plugin.Plugin
\end{lstlisting}

\subsection{\texttt{UtilsPlugin.java}}

\begin{lstlisting}[language=java]
package com.ejemplo;

import com.retroeditor.plugin.Plugin;
import com.retroeditor.plugin.PluginContext;
import javafx.application.Platform;
import javafx.scene.control.*;

public class UtilsPlugin implements Plugin {

    private PluginContext ctx;

    @Override
    public void onLoad(PluginContext context) {
        this.ctx = context;
    }

    @Override
    public void onEnable() {
        Platform.runLater(() -> {
            Menu menuUtils = new Menu("Utilidades");

            MenuItem infoItem = new MenuItem("Informacion del entorno");
            infoItem.setOnAction(e -> {
                String info = "Java: "    + System.getProperty("java.version")
                            + "\nOS: "   + System.getProperty("os.name")
                            + "\nDir: "  + System.getProperty("user.dir");
                ctx.log(info);

                Alert alert = new Alert(Alert.AlertType.INFORMATION, info, ButtonType.OK);
                alert.setTitle("Entorno");
                alert.initOwner(ctx.getPrimaryStage());
                alert.showAndWait();
            });

            menuUtils.getItems().add(infoItem);
            ctx.registerMenu(menuUtils);
        });
    }

    @Override
    public void onDisable() {
        ctx.log("UtilsPlugin desactivado.");
    }

    @Override
    public String getName() {
        return "Plugin de Utilidades";
    }
}
\end{lstlisting}

\subsection{\texttt{META-INF/services/com.retroeditor.plugin.Plugin}}

\begin{lstlisting}
com.ejemplo.UtilsPlugin
\end{lstlisting}

\subsection{Compilar e instalar}

\begin{lstlisting}[language=bash]
cd plugin-utilidades
mvn package
copy target\plugin-utilidades-1.0.0.jar ..\RetroEditor\plugins\
\end{lstlisting}

\clearpage

\section{Buenas prácticas y limitaciones}

\subsection{Buenas prácticas}

\begin{itemize}
\item \textbf{Siempre usa \texttt{Platform.runLater()}} para cualquier modificación de la UI de JavaFX desde métodos del ciclo de vida del plugin.
\item \textbf{Limpia tus recursos en \texttt{onDisable()}}: cierra hilos, streams, conexiones de red, etc.
\item \textbf{Usa \texttt{context.registerMenu()}} en lugar de acceder directamente a \texttt{getMenuBar() \\ .getMenus().add()} para que SamaruC pueda limpiar automáticamente al desactivar el plugin.
\item \textbf{Declara \texttt{provided}} todas las dependencias que ya provee el host (JavaFX, SamaruC, etc.).
\item \textbf{Versiona tu plugin} y sigue convenciones de nombre de paquete para evitar conflictos de clases.
\end{itemize}

\subsection{Limitaciones actuales}

\begin{itemize}
\item Los plugins \textbf{no pueden reemplazar} el editor de código ni el árbol de proyecto (no existen puntos de extensión para esos componentes todavía).
\item No hay sistema de dependencias entre plugins: si tu plugin depende de otro, ambos JARs deben estar presentes y el orden de carga no está garantizado.
\item El classloader usa estrategia \textbf{child-first}, por lo que puedes incluir versiones propias de librerías en tu JAR; sin embargo, las clases de \texttt{java.*}, \texttt{javax.*}, \texttt{javafx.*} y \texttt{org.kordamp.*} siempre se cargan desde el host.
\item Los plugins se cargan \textbf{desde el directorio de trabajo} (\texttt{plugins/} junto al ejecutable), no desde rutas arbitrarias del sistema.
\item Actualmente no existe una API para acceder al contenido del editor de texto abierto. Esto se planea para futuras versiones.
\end{itemize}

\vspace{2cm}
\begin{center}

\end{center}

\end{document}